% from ega4-21b-fr.tex
\oldpage[IV-4]{298}
$A$-algèbre qui est un $A$-module libre de rang fini;
en outre, on peut supposer que $D'=\operatorname{div}(t')$, où $t'$ est un élément régulier de $A'$;
on a alors $f_*(D')=\operatorname{div}(t)$, où $t=\mathrm{N}_{A'/A}(t')$ est un élément régulier de $A$ \sref{IV.21.5.2}.
On peut se borner au cas où $t'$ n'est pas inversible dans $A'$, ce qui équivaut au fait que $t$ ne l'est pas dans $A$;
l'anneau $A/tA$ est alors de longueur finie et $\neq0$, et $A'/tA'$ est composé direct d'anneaux locaux artiniens $A'_i$ $(1\leq i\leq r)$, le corps résiduel de $A'_i$ étant $\boldsymbol{k}(x'_i)$, où les $x'_i$ $(1\leq i\leq r)$ sont les points de $X'$ au-dessus de $x$.
Si $t'_i$ $(1\leq i\leq r)$ est l'image de $t'$ dans $A'_i$, $A'/t'A'$ est composé direct des $A'_i/t'_iA'_i$;
comme le produit $\operatorname{long}_{A'_i}(A'_i/t'_iA'_i)\cdot[\boldsymbol{k}(x'_i):\boldsymbol{k}(x)]$ est égal à $\operatorname{long}_A(A'_i/t'_iA'_i)$, on voit que la multiplicité au point $x$ du premier membre de \sref{IV.21.10.17.1} est $\operatorname{long}_A(A'/t'A')$, si bien que la formule à démontrer se réduit à
\[
\label{IV.21.10.17.2}
  \operatorname{long}_A(A'/t'A')=\operatorname{long}_A(A/tA).
  \tag{21.10.17.2}
\]

Cette relation découle du lemme plus général suivant.
\end{proof}

\begin{lemma}[21.10.17.3]
\label{IV.21.10.17.3}
Soient $A$ un anneau local noethérien de dimension $1$, $M$ un $A$-module libre de rang fini, $u$ un endomorphisme injectif de $M$.
Alors on a
\[
\label{IV.21.10.17.4}
  \operatorname{long}_A(\operatorname{Coker}u)=\operatorname{long}_A(A/(\det u)A).
  \tag{21.10.17.4}
\]
\end{lemma}

\begin{proof}
Distinguons plusieurs cas.

I) $A$ est un anneau de valuation discrète.
Soit en effet $\pi$ une uniformisante de $A$, et remarquons que $\operatorname{long}_A(A/\pi^kA)=k$;
si $n$ est le rang de $M$ et si $\pi^{m_i}$ $(1\leq i\leq n)$ sont les facteurs invariants de $u$, $\operatorname{Coker}(u)$ est somme directe des $A$-modules $A/\pi^{m_i}A$, donc a pour longueur $m=\sum_{i=1}^nm_i$ et $\det(u)$ est produit d'un élément inversible et de $\pi^m$, d'où la conclusion dans ce cas (Bourbaki, \emph{Alg.}, chap.~VII, \textsection~4, n\textsuperscript{o}~5, cor.~1 de la prop.~4).

II) $A$ est un anneau complet intègre (de dimension $1$).
On sait alors $(\mathbf{0},19.8.8)$, (ii) qu'il y a un sous-anneau $B$ de $A$, qui est un anneau de valuation discrète, tel que $B\to A$ soit un homomorphisme local faisant de $A$ un $B$-module de type fini;
comme ce $B$-module est évidemment sans torsion, il est libre (Bourbaki, \emph{Alg. comm.}, chap.~VI, \textsection~3, n\textsuperscript{o}~6, lemme~1).
Notons $M'$ l'ensemble $M$ muni de sa structure de $B$-module (libre), par $u'$ l'endomorphisme $u$ considéré comme $B$-endomorphisme.
Il résulte de I) que l'on a
\[
\label{IV.21.10.17.5}
  \operatorname{long}_B(\operatorname{Coker}u')=\operatorname{long}_B(B/(\det u')B).
  \tag{21.10.17.5}
\]
Mais on a (Bourbaki, \emph{Alg.}, chap.~VIII, \textsection~12, n\textsuperscript{o}~2, prop.~7)
\[
  \det u'=\mathrm{N}_{A/B}(\det u),
\]
donc, en appliquant \sref{IV.21.10.17.4} à l'homothétie $x\mapsto(\det u)x$ du $B$-module libre $A$, il vient
\[
  \operatorname{long}_B(B/(\det u')B)=\operatorname{long}_B(A/(\det u)A),
\]
d'où, en portant dans \sref{IV.21.10.17.5} et divisant par $[\boldsymbol{k}(A):\boldsymbol{k}(B)]$, longueur du $B$-module $\boldsymbol{k}(A)$, on obtient \sref{IV.21.10.17.4} dans le cas envisagé.

